% Volume 1: Quantity
% Opener: tagline + scope. Sits immediately after the \part{} declaration.

\vspace*{1.5cm}
\begin{center}
\itshape\large
Reality is measurable
\end{center}
\vspace{0.5cm}

\noindent
We teach the reader to convert reality into numbers and to know what those numbers mean. Measurement, units, error, calibration, dimensional analysis, and the statistics every working engineer needs.

\vspace{0.5cm}

\section*{What this volume is for}

Engineering begins with measurement. We open the work here, in the
volume that turns the world into numbers, because everything in the
remaining eleven volumes assumes a reader who already carries a
measurement habit. Mathematics, force, energy, matter, life,
information, machines, systems, failure, design, and civilization are
each a discipline that operates on quantities. A reader who arrives
at Volume II without a working sense of what a measurement is, what
its uncertainty means, and what it costs to communicate, will read
the rest of the book as theorem followed by theorem. A reader who
spends time here will read the rest as a sequence of decisions about
real, dated, error-bearing numbers. The latter is what we want.

Volume I is the smallest of the twelve. We have no encyclopedia of
metrology to deliver. We have a set of habits to install: estimation
before calculation, dimensions before equations, calibration before
trust, error as a first-class quantity carried alongside every
result, statistics that the reader can defend rather than perform,
and the discipline of an engineering notebook that an outside reader
could pick up and use. The volume is short because the work it asks
of the reader is hard, and we prefer to ask the work in nine
focused chapters rather than thirty padded ones.

\section*{What the reader brings}

The reader who opens this book is assumed to have an adult life and
adult patience, secondary-school arithmetic, geometry, trigonometry,
and an ordinary sense of physical quantities of the kind a literate
adult acquires by living in the world. We do not assume calculus,
linear algebra, statistics, or programming. Volume II takes those up
in order. The reader is assumed to have, somewhere in their life, a
kitchen scale, a thermometer, a tape measure, a stopwatch, and a
phone. None of those is required to read the chapters, but the
volume project asks the reader to use them, and the projects in the
later volumes assume the habits this one builds.

The reader is also assumed to bring an existing intuition for
quantities the world hands them every day: hot and cold, heavy and
light, fast and slow, near and far, loud and quiet. Engineering does
not replace those intuitions. Engineering refines them by attaching
them to numbers, units, and uncertainty, and then asks the reader
which of the original intuitions were correct, which were biased,
and which were missing entirely. By the end of the volume, every
intuition the reader brought in will have been touched at least
once.

\section*{The arc of the nine chapters}

Chapter 1 establishes that engineering and measurement are the same
discipline at different scales. Three named accidents make the
argument vivid: the Mars Climate Orbiter, the Air Canada Gimli
Glider incident, and the Ariane 5 maiden flight. Each is a story
about a number that was not measured at the interface where it
mattered.

Chapter 2 introduces units and dimensions as the syntax of
quantitative reasoning. The reader meets the SI base units, the
derived units, the practice of carrying units through every step of
every calculation, and dimensional analysis as a debugging tool. The
Buckingham pi theorem appears here, deliberately early, because its
power as a sanity check shows up immediately and its formal
mathematical content is light.

Chapter 3 introduces calibration and traceability: how a measurement
gets its meaning. The chapter walks the chain from the laboratory
instrument through the working standard, the secondary standard, the
primary standard, and finally the international definitions of the
SI units themselves. Calibration is treated as the act that links a
local reading to the global community of measurement.

Chapter 4 introduces error and uncertainty. The reader learns the
difference between the two, the propagation rules for sums,
products, and general functions, the distinction between Type A and
Type B uncertainty per the international guide, and the skeleton of
an uncertainty budget. This is the most demanding chapter of the
volume, and the most consequential.

Chapters 5 through 7 give the reader concrete instruments and the
quantities they touch. Chapter 5 covers sensors and instruments
(thermometers, pressure gauges, load cells, photodiodes, microphones,
inertial sensors, the smartphone as instrument cluster). Chapter 6
covers time, frequency, and signals: the clock as the first
instrument, sampling, the difference between an analogue and a
digital reading. Chapter 7 covers the geometry of measurement:
length, area, volume, and mass, with their corresponding instruments
and their characteristic failure modes.

Chapter 8 introduces statistics. The standard is engineering
statistics: enough to summarise, infer, and decide, not enough to
write a graduate paper in statistical theory. Distributions, the
mean and its uncertainty, hypothesis testing as a tool the engineer
uses sparingly, and a brief encounter with Bayesian reasoning where
prior information is available.

Chapter 9 closes the volume by turning estimation into a discipline.
The reader has done estimation throughout the previous chapters; in
the closing chapter we make estimation explicit, give it a vocabulary,
and exercise it on problems that range across all twelve future
volumes. By the end of Chapter 9 the reader can produce an
order-of-magnitude answer, with a confidence range, to most
quantitative questions an engineer is likely to face in their first
years of practice.

\section*{What the reader will be able to do at the end}

A reader who finishes Volume I should be able, without thinking
about it, to:

\begin{itemize}[leftmargin=*]
  \item Estimate any engineering quantity to within a factor of three
    before reaching for a calculator.
  \item Carry units through every step of every calculation, and
    detect a unit error before the calculation finishes.
  \item Report any measured quantity with its uncertainty and
    coverage factor, and describe how the uncertainty was derived.
  \item Read an instrument's specification sheet and identify the
    quantities the instrument measures, the quantities it does not,
    and the conditions under which its specification holds.
  \item Distinguish precision from accuracy and apply the
    distinction to a real-world reading.
  \item Trace a calibration through a documented chain to a
    primary standard.
  \item Keep a measurement notebook of a kind that an outside
    engineer could pick up and use without ambiguity.
  \item Recognise, in the field, the small daily form of the
    measurement failures that the named accidents in Chapter 1
    illustrate at scale.
\end{itemize}

A reader who reaches Volume II without these habits has read Volume I
without doing it. We say so plainly because the rest of the book
will assume otherwise. The volume project, described below, is the
mechanism by which the reader proves to themselves that the habits
have taken.

\section*{Archetypes introduced in this volume}

Three of the book's ten recurring problem archetypes make their first
appearance here. They are introduced as formal environments,
\verb|\begin{archetype}|, and they will recur in every later volume.

\begin{itemize}[leftmargin=*]
  \item \engterm{Scaling}: orders of magnitude, dimensional reasoning,
    and the engineer's internal ruler. Introduced in Chapter 1 and
    Chapter 2; revisited in Chapter 9 as the basis of estimation.
  \item \engterm{Failure}: the ways measurement itself fails, and the
    silent unit confusion at interfaces that is the discipline's
    most common chronic accident. Introduced in Chapter 1 and
    extended in Chapter 4 as the propagation of uncertainty.
  \item \engterm{Uncertainty}: the parameter that turns a number into
    a measurement, the engineer's first acknowledgement of
    imperfect knowledge. Introduced in Chapter 4 and revisited in
    Chapter 8 as statistical inference.
\end{itemize}

The remaining seven archetypes (balance, transport, stability,
optimisation, control, interface, ethics) appear from Volume II
onwards. The reader can read this list with the comfortable
knowledge that the book's conceptual machinery will arrive
gradually, in the chapters where the machinery does the work.

\section*{The volume project}

A reader's measurement notebook accompanies the volume. By the close
of Chapter 9 the notebook should contain, at minimum: one calibration
of a household instrument; one error-bounded measurement of a
non-trivial quantity; three order-of-magnitude estimates compared to
a published or measured ground truth; and one fully documented
experiment with at least thirty trials and a proper statistical
analysis. The notebook is graded against a public rubric. The general
editor's reference notebook is published online; the reader is
encouraged to compare it to their own only after completing the
project, never before. The notebook is the volume's deliverable, and
the credential by which the reader earns admission to Volume II.

\section*{A note on the next volume}

Volume II takes up the mathematics that this volume's quantities
demand. Calculus, linear algebra, probability, optimisation, and
numerical methods. The bridge is direct: every quantity introduced
here will, in Volume II, be operated on by a function, organised
into a vector or a matrix, drawn from a distribution, or computed by
an algorithm. The transition is from \engterm{we have numbers} to
\engterm{we have functions of numbers}. The reader who carries the
habits of this volume across that bridge has already done the
hardest part of becoming an engineer.

\vspace{1cm}
